% ==============================================================================
% Figura: Diagrama de Bloques Atómico — Estado FURUTA_STATE_BALANCE_PID (1 kHz)
% Estrategia M4 Ultra-Silk: 1 función = 1 bloque atómico, solo nombre corto.
% Fiel a furuta_control.c, as5600.c, pid_controller.c y l298n_motor.c.
% Archivo modular listo para \input{} en IEEEtran dos columnas (figure* [!t]).
% Bibliotecas TikZ requeridas: arrows.meta, positioning, fit, calc,
%                            shapes.geometric, backgrounds, babel, bending.
% ==============================================================================
\begin{figure*}[!t]
\centering
% --- Paleta semántica armónica (local al entorno figure*) ---
\definecolor{blkInk}{HTML}{1E293B}
\definecolor{blkPendEdge}{HTML}{1B4965}
\definecolor{blkPendFill}{HTML}{EAF4FB}
\definecolor{blkArmEdge}{HTML}{0F766E}
\definecolor{blkArmFill}{HTML}{E6F7F5}
\definecolor{blkCondEdge}{HTML}{B45309}
\definecolor{blkCondFill}{HTML}{FEF6EB}
\definecolor{blkDistEdge}{HTML}{7E2268}
\definecolor{blkDistFill}{HTML}{FBF0F8}
\definecolor{blkActEdge}{HTML}{A63429}
\definecolor{blkActFill}{HTML}{FDF0EE}
\definecolor{blkPlantEdge}{HTML}{243447}
\definecolor{blkPlantFill}{HTML}{ECEFF4}
\definecolor{blkZoneZ1}{HTML}{33658A}
\definecolor{blkZoneL2}{HTML}{A63429}
\definecolor{blkZoneL3}{HTML}{2A6B43}

% --- Parámetros centralizados de ritmo geométrico (Layout by Construction) ---
\def\gapErrX{0.76cm}
\def\gapLtwoX{0.80cm}
\def\gapTapX{1.16cm}
\def\zOneRowGap{1.12cm}
\def\boxH{0.88cm}
\def\subH{0.76cm}
\def\sumSz{0.56cm}
\def\hopRad{0.085cm}
\def\trackStepX{0.34cm}
\def\refLeadX{1.48cm}

\resizebox{\textwidth}{!}{%
\begin{tikzpicture}[
    >=Stealth,
    font=\sffamily,
    % --- Estilos de bloques atómicos (1 función = 1 bloque, solo nombre corto) ---
    baseblk/.style={
        rectangle,
        rounded corners=3.5pt,
        align=center,
        inner xsep=3.5pt,
        inner ysep=3pt,
        text=blkInk,
        font=\normalsize\sffamily\bfseries,
        minimum height=\boxH
    },
    pendblk/.style={baseblk, draw=blkPendEdge, line width=1.1pt, fill=blkPendFill, text width=1.74cm},
    armblk/.style={baseblk, draw=blkArmEdge, line width=1.1pt, fill=blkArmFill, text width=1.60cm},
    condblk/.style={baseblk, draw=blkCondEdge, line width=1.1pt, fill=blkCondFill, text width=1.62cm},
    distblk/.style={baseblk, draw=blkDistEdge, line width=1.05pt, fill=blkDistFill, text width=2.15cm, minimum height=\subH},
    actblk/.style={baseblk, draw=blkActEdge, line width=1.1pt, fill=blkActFill, text width=2.04cm},
    plantblk/.style={baseblk, draw=blkPlantEdge, line width=1.35pt, fill=blkPlantFill, text width=1.98cm},
    sensTblk/.style={baseblk, draw=blkArmEdge, line width=1.1pt, fill=blkArmFill, text width=2.36cm},
    sensAblk/.style={baseblk, draw=blkPendEdge, line width=1.1pt, fill=blkPendFill, text width=2.36cm},
    sum/.style={
        circle,
        draw=blkInk,
        line width=1.05pt,
        fill=white,
        minimum size=\sumSz,
        inner sep=0pt,
        font=\small\bfseries,
        text=blkInk
    },
    dot/.style={circle, fill=#1, inner sep=1.5pt},
    % --- Flechas canónicas ---
    sig/.style={-{Stealth[length=4.8pt, width=3.5pt]}, line width=1.0pt, draw=blkInk},
    sigshort/.style={-{Stealth[length=4.0pt, width=3.0pt]}, line width=1.0pt, draw=blkInk},
    fbA/.style={-{Stealth[length=4.4pt, width=3.2pt]}, line width=0.95pt, draw=blkPendEdge},
    fbT/.style={-{Stealth[length=4.4pt, width=3.2pt]}, line width=0.95pt, draw=blkArmEdge},
    lineA/.style={line width=0.95pt, draw=blkPendEdge},
    lineT/.style={line width=0.95pt, draw=blkArmEdge},
    adapt/.style={-{Stealth[length=4.2pt, width=3.0pt]}, line width=0.95pt, densely dashed, draw=blkDistEdge},
    % --- Contenedores y etiquetas ---
    zonebox/.style={
        rectangle,
        rounded corners=5pt,
        draw=#1,
        densely dashed,
        line width=0.95pt,
        fill=#1,
        fill opacity=0.035,
        text opacity=1
    },
    zonetitle/.style={
        rectangle,
        rounded corners=3pt,
        draw=#1,
        line width=0.8pt,
        fill=white,
        inner xsep=5.5pt,
        inner ysep=2.2pt,
        font=\small\sffamily\bfseries,
        text=#1
    },
    signlbl/.style={font=\footnotesize\bfseries, text=blkInk, inner sep=1pt},
    wirelbl/.style={font=\small\sffamily, text=blkInk, inner sep=1.5pt}
]

% ==============================================================================
% NIVEL 1: BALANCE DUAL PARALELO + ACONDICIONAMIENTO ADAPTATIVO — ÚNICO ANCLAJE (0,0)
% ==============================================================================
\node[sum] (sig1) at (0,0) {$\Sigma_1$};
\node[pendblk, right=\gapErrX of sig1] (pd_pend) {PD\\P\'endulo};

\node[armblk, below=\zOneRowGap of pd_pend] (pid_arm) {PID\\Brazo};
\node[sum] (sig2) at (sig1 |- pid_arm) {$\Sigma_2$};
\node[armblk, right=0.50cm of pid_arm] (sat_arm) {Sat.\\Brazo};

\coordinate (z1_mid_y) at ($(pd_pend.east)!0.5!(pid_arm.east)$);
\node[sum, right=0.72cm of sat_arm.east |- z1_mid_y] (sig_u) {$\Sigma_u$};

\node[condblk, right=0.78cm of sig_u] (notch1) {Notch 1};
\coordinate[right=0.42cm of notch1.east] (n_split);

\node[condblk, anchor=west] (notch2) at ($(n_split |- pid_arm)+(0.36cm, 0)$) {Notch 2};
\node[condblk, anchor=west] (mezcla) at ($(notch2.east |- z1_mid_y)+(0.82cm, 0)$) {Mezcla};
\node[condblk, right=0.60cm of mezcla] (sat_adap) {Sat.\\Adap.};
\node[condblk, right=0.54cm of sat_adap] (slew_adap) {Slew\\Adap.};

\coordinate (dist_bus_y) at ($(mezcla.south)+(0, -0.45cm)$);
\node[distblk, below=0.74cm of notch2.south -| mezcla.center] (dist_blk) {\'Indice dist};

\coordinate (z1_top_pad)   at ([yshift=0.92cm]pd_pend.north);
\coordinate (z1_bot_pad)   at ([yshift=-0.26cm]dist_blk.south);
\coordinate (z1_left_pad)  at ([xshift=-0.38cm]sig1.west);
\coordinate (z1_right_pad) at ([xshift=0.30cm]slew_adap.east);

\begin{scope}[on background layer]
    \node[zonebox=blkZoneZ1, inner sep=0pt,
          fit=(z1_top_pad)(z1_bot_pad)(z1_left_pad)(z1_right_pad)] (zone1) {};
\end{scope}
\node[zonetitle=blkZoneZ1, anchor=west] at ([xshift=0.35cm]zone1.north west)
    {Balance (Paralelo + Cascada)};

% ==============================================================================
% NIVEL 2: ACTUACIÓN Y PLANTA EN CASCADA
% ==============================================================================
\coordinate (l2_y) at ([yshift=-2.00cm]zone1.south);
\node[actblk, anchor=west] (banda) at (pd_pend.west |- l2_y) {Banda\\Silenciosa};
\node[actblk, right=\gapLtwoX of banda] (rampa) {Rampa\\L298N};
\node[actblk, right=\gapLtwoX of rampa] (pwm) {Modulador\\PWM};
\node[actblk, right=\gapLtwoX of pwm] (puente) {Puente H\\L298N};
\node[actblk, right=\gapLtwoX of puente] (motor) {Motor\\80:1};
\node[plantblk, right=\gapLtwoX of motor] (planta) {P\'endulo\\Furuta};

\coordinate (l2_top_pad)   at ([yshift=0.36cm]banda.north);
\coordinate (l2_bot_pad)   at ([yshift=-0.28cm]banda.south);
\coordinate (l2_left_pad)  at ([xshift=-0.28cm]banda.west);
\coordinate (l2_right_pad) at ([xshift=0.28cm]planta.east);

\begin{scope}[on background layer]
    \node[zonebox=blkZoneL2, inner sep=0pt,
          fit=(l2_top_pad)(l2_bot_pad)(l2_left_pad)(l2_right_pad)] (zoneL2) {};
\end{scope}
\node[zonetitle=blkZoneL2, anchor=west] at ([xshift=0.32cm]zoneL2.north west)
    {Actuaci\'on y Planta (Cascada)};

% ==============================================================================
% NIVEL 3: REALIMENTACIÓN Y FILTRADO (DERECHA -> IZQUIERDA)
% ==============================================================================
\coordinate (l3_t_y) at ([yshift=-1.95cm]zoneL2.south);
\coordinate (as5600_east_x) at ([xshift=-1.00cm]planta.south);

\node[sensTblk, anchor=east] (as5600_base) at (as5600_east_x |- l3_t_y) {AS5600\\Base};
\node[sensTblk, left=\gapTapX of as5600_base] (desenv) {Desenvoltura};
\node[sensTblk, left=\gapTapX of desenv] (deriv_t) {Derivador $\theta$};
\node[sensTblk, left=\gapTapX of deriv_t] (lpf1_t) {LPF IIR 1\\($\theta$)};
\node[sensTblk, left=\gapTapX of lpf1_t] (lpf2_t) {LPF IIR 2\\($\theta$)};

\node[sensAblk, below=0.80cm of as5600_base] (as5600_dist) {AS5600\\Distal};
\node[sensAblk, left=\gapTapX of as5600_dist] (envolt) {Envoltura};
\node[sensAblk, left=\gapTapX of envolt] (deriv_a) {Derivador $\alpha$};
\node[sensAblk, left=\gapTapX of deriv_a] (lpf1_a) {LPF IIR 1\\($\alpha$)};
\node[sensAblk, left=\gapTapX of lpf1_a] (lpf2_a) {LPF IIR 2\\($\alpha$)};

\coordinate (l3_top_pad)   at ([yshift=0.72cm]lpf2_t.north);
\coordinate (l3_bot_pad)   at ([yshift=-1.08cm]lpf2_a.south);
\coordinate (l3_left_pad)  at ([xshift=-0.96cm]lpf2_t.west);
\coordinate (l3_right_pad) at ([xshift=0.24cm]as5600_base.east);

\begin{scope}[on background layer]
    \node[zonebox=blkZoneL3, inner sep=0pt,
          fit=(l3_top_pad)(l3_bot_pad)(l3_left_pad)(l3_right_pad)] (zoneL3) {};
\end{scope}
\node[zonetitle=blkZoneL3, anchor=west] at ([xshift=0.35cm]zoneL3.south west)
    {Realimentaci\'on y Filtrado (Paralelo + Cascada)};

% ==============================================================================
% CONEXIONES DE NIVEL 1 (CONTROL DUAL PARALELO + DUAL-NOTCH ADAPTATIVO)
% ==============================================================================
\coordinate (ref_start_x) at ([xshift=-\refLeadX]sig1.west);

% Entradas de referencia
\draw[sig] (ref_start_x |- sig1) -- (sig1.west)
    node[wirelbl, pos=0.38, above=2pt] {$\tilde{\alpha}_{\mathrm{ref}}$};
\node[signlbl, anchor=south] at ([xshift=-0.13cm, yshift=0.06cm]sig1.west) {$+$};

\draw[sig] (ref_start_x |- sig2) -- (sig2.west)
    node[wirelbl, pos=0.38, above=2pt] {$\theta_{\mathrm{ref}}$};
\node[signlbl, anchor=south] at ([xshift=-0.13cm, yshift=0.06cm]sig2.west) {$+$};

% Errores hacia controladores
\draw[sigshort] (sig1.east) -- (pd_pend.west)
    node[wirelbl, midway, above=2pt] {$e_\alpha$};
\draw[sigshort] (sig2.east) -- (pid_arm.west)
    node[wirelbl, midway, above=2pt] {$e_\theta$};
\draw[sigshort] (pid_arm.east) -- (sat_arm.west);

% Convergencia paralela en Sigma_u
\draw[sigshort, rounded corners=4pt] (pd_pend.east) -| (sig_u.north)
    node[wirelbl, pos=0.28, above=2pt] {$u_\alpha$};
\node[signlbl, anchor=south west] at ([xshift=0.08cm, yshift=0.06cm]sig_u.north) {$+$};

\draw[sigshort, rounded corners=4pt] (sat_arm.east) -| (sig_u.south)
    node[wirelbl, pos=0.28, below=2pt] {$u_\theta$};
\node[signlbl, anchor=north west] at ([xshift=0.08cm, yshift=-0.06cm]sig_u.south) {$+$};

\draw[sigshort] (sig_u.east) -- (notch1.west)
    node[wirelbl, midway, above=2pt] {$u_{\mathrm{raw}}$};

% Bifurcación paralela Dual-Notch en n_split -> Mezcla
\draw[line width=1.0pt, draw=blkInk] (notch1.east) -- (n_split);
\node[dot=blkInk] at (n_split) {};

\coordinate (mez_in_top) at ([yshift= 0.16cm]mezcla.west);
\coordinate (mez_in_bot) at ([yshift=-0.16cm]mezcla.west);

\draw[sigshort, rounded corners=4pt]
    (n_split) -- (n_split |- mez_in_top) -- (mez_in_top)
    node[wirelbl, pos=0.45, above=2pt] {$u_{n11}$};

\draw[sigshort, rounded corners=4pt]
    (n_split) |- (notch2.west);

\draw[sigshort, rounded corners=4pt]
    (notch2.east) -| ([xshift=-0.42cm]mez_in_bot) -- (mez_in_bot);
\node[wirelbl, anchor=north]
    at ($(notch2.east)+(0.20cm, -0.05cm)$) {$u_{n6}$};

\draw[sigshort] (mezcla.east) -- (sat_adap.west)
    node[wirelbl, midway, above=2pt] {$u_{\mathrm{sh}}$};
\draw[sigshort] (sat_adap.east) -- (slew_adap.west);

% Bus adaptativo punteado dist desde Índice dist
\coordinate (dist_rail_l) at (dist_blk.north |- dist_bus_y);
\coordinate (dist_rail_r) at (slew_adap.south |- dist_bus_y);
\draw[line width=0.95pt, densely dashed, draw=blkDistEdge]
    (dist_blk.north) -- (dist_rail_l) -- (dist_rail_r);
\node[wirelbl, text=blkDistEdge, anchor=east]
    at ([xshift=-0.08cm]$(dist_blk.north)!0.48!(dist_rail_l)$) {$\mathrm{dist}$};
\foreach \tgt in {mezcla, sat_adap} {
    \node[dot=blkDistEdge] at (\tgt.south |- dist_bus_y) {};
    \draw[adapt] (\tgt.south |- dist_bus_y) -- (\tgt.south);
}
\draw[adapt] (dist_rail_r) -- (slew_adap.south);

% ==============================================================================
% CONEXIONES DE NIVEL 2 (u[k] -> ACTUACIÓN L298N -> PLANTA FURUTA)
% ==============================================================================
\coordinate (u_right_x) at ([xshift=0.62cm]slew_adap.east);
\coordinate (lane_u_y)  at ([yshift=0.58cm]zoneL2.north);
\coordinate (u_in_x)    at ([xshift=-0.52cm]banda.west);

\draw[sig, rounded corners=4pt]
    (slew_adap.east) -- (u_right_x |- slew_adap.east)
    -- (u_right_x |- lane_u_y)
    -- (u_in_x |- lane_u_y)
    |- (banda.west);
\node[wirelbl, anchor=south]
    at ($(motor.east |- lane_u_y)!0.50!(planta.east |- lane_u_y) + (0, 0.04cm)$) {$u[k]$};

\draw[sigshort] (banda.east) -- (rampa.west);
\draw[sigshort] (rampa.east) -- (pwm.west);
\draw[sigshort] (pwm.east) -- (puente.west);
\draw[sigshort] (puente.east) -- (motor.west)
    node[wirelbl, midway, above=2pt] {$V_m$};
\draw[sigshort] (motor.east) -- (planta.west)
    node[wirelbl, midway, above=2pt] {$\tau_0$};

% ==============================================================================
% CONEXIONES DE NIVEL 3 (PLANTA -> ENCODERS -> FILTROS -> 5 LÍNEAS EXPLÍCITAS)
% ==============================================================================
\coordinate (pl_t_out) at ([xshift=-0.28cm]planta.south);
\coordinate (pl_a_out) at ([xshift= 0.28cm]planta.south);

\draw[fbT, rounded corners=4pt] (pl_t_out) |- (as5600_base.east);
\node[wirelbl, text=blkArmEdge, anchor=east]
    at ([xshift=-0.06cm]$(pl_t_out |- zoneL2.south)!0.48!(pl_t_out |- zoneL3.north)$) {$\theta(t)$};

\draw[fbA, rounded corners=4pt] (pl_a_out) |- (as5600_dist.east);
\node[wirelbl, text=blkPendEdge, anchor=west]
    at ([xshift=0.06cm]$(pl_a_out |- as5600_base.east)!0.50!(pl_a_out |- as5600_dist.east)$) {$\alpha(t)$};

% Rama theta (Derecha -> Izquierda)
\draw[fbT] (as5600_base.west) -- (desenv.east);
\draw[fbT] (desenv.west) -- (deriv_t.east);
\draw[fbT] (deriv_t.west) -- (lpf1_t.east);
\draw[fbT] (lpf1_t.west) -- (lpf2_t.east);

% Rama alpha (Derecha -> Izquierda)
\draw[fbA] (as5600_dist.west) -- (envolt.east);
\draw[fbA] (envolt.west) -- (deriv_a.east);
\draw[fbA] (deriv_a.west) -- (lpf1_a.east);
\draw[fbA] (lpf1_a.west) -- (lpf2_a.east);

% --- 5 pistas verticales ordenadas en el margen izquierdo (sin bus opaco) ---
\coordinate (trk_dt2_x)   at ([xshift=-0.34cm]ref_start_x);
\coordinate (trk_theta_x) at ([xshift=-\trackStepX]trk_dt2_x);
\coordinate (trk_da1_x)   at ([xshift=-\trackStepX]trk_theta_x);
\coordinate (trk_alpha_x) at ([xshift=-\trackStepX]trk_da1_x);
\coordinate (trk_da2_x)   at ([xshift=-\trackStepX]trk_alpha_x);

% Puntos de derivación en Nivel 3
\coordinate (tap_theta) at ($(desenv.west)!0.68!(deriv_t.east)$);
\coordinate (tap_alpha) at ($(envolt.west)!0.74!(deriv_a.east)$);
\coordinate (tap_da1)   at ($(lpf1_a.west)!0.68!(lpf2_a.east)$);

\coordinate (rail_theta_y) at ([yshift=0.44cm]lpf2_t.north);
\coordinate (rail_da1_y)   at ([yshift=-0.42cm]lpf2_a.south);
\coordinate (rail_alpha_y) at ([yshift=-0.78cm]lpf2_a.south);

\node[dot=blkArmEdge]  at (tap_theta) {};
\node[dot=blkPendEdge] at (tap_alpha) {};
\node[dot=blkPendEdge] at (tap_da1) {};

% Etiquetas en las salidas de Nivel 3
\node[wirelbl, text=blkArmEdge, anchor=west]
    at ([xshift=0.06cm]$(tap_theta)!0.48!(tap_theta |- rail_theta_y)$) {$\theta$};
\node[wirelbl, text=blkArmEdge, anchor=south]
    at ([xshift=-0.46cm, yshift=0.04cm]lpf2_t.west) {$\dot{\theta}_{f2}$};

\node[wirelbl, text=blkPendEdge, anchor=west]
    at ([xshift=0.06cm]$(tap_alpha)!0.28!(tap_alpha |- rail_alpha_y)$) {$\alpha,\,\tilde{\alpha}$};
\node[wirelbl, text=blkPendEdge, anchor=west]
    at ([xshift=0.06cm]$(tap_da1)!0.44!(tap_da1 |- rail_da1_y)$) {$\dot{\alpha}_{f1}$};
\node[wirelbl, text=blkPendEdge, anchor=north]
    at ([xshift=-0.46cm, yshift=-0.04cm]lpf2_a.west) {$\dot{\alpha}_{f2}$};

% --- Coordenadas Y de destino en Nivel 1 ---
\coordinate (y_pd_pend_fb) at ([yshift=0.38cm]pd_pend.north);
\coordinate (y_sig1_fb)    at ([yshift=-0.42cm]sig1.south);
\coordinate (y_sig2_fb)    at ([yshift=-0.42cm]sig2.south);
\coordinate (y_brazo_fb)   at ([yshift=-0.72cm]pid_arm.south);
\coordinate (y_dist_top)   at ([yshift= 0.18cm]dist_blk.west);
\coordinate (y_dist_bot)   at ([yshift=-0.18cm]dist_blk.west);

% 1) Pista dtheta_f2 -> PID Brazo.south (interior, termina en y_brazo_fb)
\draw[fbT, rounded corners=4pt]
    (lpf2_t.west) -- (trk_dt2_x |- lpf2_t.west)
    -- (trk_dt2_x |- y_brazo_fb) -| (pid_arm.south);
\node[wirelbl, text=blkArmEdge, anchor=south east]
    at ([xshift=-0.08cm, yshift=0.02cm]pid_arm.south |- y_brazo_fb) {$\dot{\theta}_{f2}$};

% 2) Pista theta -> Sigma_2.south (termina en y_sig2_fb, por encima de y_brazo_fb)
\draw[fbT, rounded corners=4pt]
    (tap_theta) -- (tap_theta |- rail_theta_y)
    -- (trk_theta_x |- rail_theta_y)
    -- (trk_theta_x |- y_sig2_fb) -| (sig2.south);
\node[signlbl, text=blkArmEdge, anchor=north west] at ([xshift=0.07cm, yshift=-0.06cm]sig2.south) {$-$};
\node[wirelbl, text=blkArmEdge, anchor=south]
    at ($(trk_dt2_x |- y_sig2_fb)!0.56!(sig2.south |- y_sig2_fb) + (0, 0.02cm)$) {$\theta$};

% 3) Pista dalpha_f1 -> Índice dist (termina en y_dist_bot, con 2 puentes sobre pistas theta)
\draw[lineA, rounded corners=4pt]
    (tap_da1) -- (tap_da1 |- rail_da1_y)
    -- (trk_da1_x |- rail_da1_y)
    -- (trk_da1_x |- y_dist_bot)
    -- ($(trk_theta_x |- y_dist_bot)+(-\hopRad,0)$)
    arc[start angle=180, end angle=0, radius=\hopRad]
    -- ($(trk_dt2_x |- y_dist_bot)+(-\hopRad,0)$)
    arc[start angle=180, end angle=0, radius=\hopRad];
\draw[fbA]
    ($(trk_dt2_x |- y_dist_bot)+(\hopRad,0)$) -- (y_dist_bot);
\node[wirelbl, text=blkPendEdge, anchor=north east]
    at ([xshift=-0.10cm, yshift=-0.02cm]y_dist_bot) {$\dot{\alpha}_{f1}$};

% 4) Pista alpha_tilde -> Índice dist (derivación con 2 puentes) y -> Sigma_1.south
\draw[lineA, rounded corners=4pt]
    (tap_alpha) -- (tap_alpha |- rail_alpha_y)
    -- (trk_alpha_x |- rail_alpha_y)
    -- (trk_alpha_x |- y_sig1_fb);
\node[dot=blkPendEdge] at (trk_alpha_x |- y_dist_top) {};

\draw[lineA]
    (trk_alpha_x |- y_dist_top)
    -- ($(trk_theta_x |- y_dist_top)+(-\hopRad,0)$)
    arc[start angle=180, end angle=0, radius=\hopRad]
    -- ($(trk_dt2_x |- y_dist_top)+(-\hopRad,0)$)
    arc[start angle=180, end angle=0, radius=\hopRad];
\draw[fbA]
    ($(trk_dt2_x |- y_dist_top)+(\hopRad,0)$) -- (y_dist_top);
\node[wirelbl, text=blkPendEdge, anchor=south east]
    at ([xshift=-0.10cm, yshift=0.02cm]y_dist_top) {$\tilde{\alpha}$};

\draw[fbA, rounded corners=4pt]
    (trk_alpha_x |- y_sig1_fb) -| (sig1.south);
\node[signlbl, text=blkPendEdge, anchor=north west] at ([xshift=0.07cm, yshift=-0.06cm]sig1.south) {$-$};
\node[wirelbl, text=blkPendEdge, anchor=south]
    at ($(trk_dt2_x |- y_sig1_fb)!0.56!(sig1.south |- y_sig1_fb) + (0, 0.02cm)$) {$\tilde{\alpha}$};

% 5) Pista dalpha_f2 -> PD Péndulo.north (exterior, con 2 puentes en Nivel 3)
\draw[lineA]
    (lpf2_a.west)
    -- ($(trk_da1_x |- lpf2_a.west)+(\hopRad,0)$)
    arc[start angle=0, end angle=180, radius=\hopRad]
    -- ($(trk_alpha_x |- lpf2_a.west)+(\hopRad,0)$)
    arc[start angle=0, end angle=180, radius=\hopRad];
\draw[fbA, rounded corners=4pt]
    ($(trk_alpha_x |- lpf2_a.west)+(-\hopRad,0)$) -- (trk_da2_x |- lpf2_a.west)
    -- (trk_da2_x |- y_pd_pend_fb) -| (pd_pend.north);
\node[wirelbl, text=blkPendEdge, anchor=west]
    at ([xshift=0.08cm]$(pd_pend.north)!0.50!(pd_pend.north |- y_pd_pend_fb)$) {$\dot{\alpha}_{f2}$};

\end{tikzpicture}%
}
\caption{Diagrama de bloques del controlador de equilibrio \texttt{BALANCE\_PID}.}
\label{fig:diagrama_bloques_control}
\end{figure*}